\documentclass[../../main.tex]{subfiles}
\begin{document}
\section{嵌入子流形}\label{section:ch05-嵌入子流形}

设 $M$ 是带边或不带边的光滑流形。$M$ 的一个\textbf{嵌入子流形}是 $M$ 的子集 $S\subseteq M$,它在子空间拓扑下是一个(不带边的)流形,并赋予一个光滑结构使包含映射 $S\hookrightarrow M$ 成为光滑嵌入。嵌入子流形也被某些作者称为正则子流形。

若 $S$ 是 $M$ 的嵌入子流形,则差 $\dim M-\dim S$ 称为 $S$ 在 $M$ 中的\textbf{余维数},包含流形 $M$ 称为 $S$ 的\textbf{外围流形}。余维数为 $1$ 的嵌入子流形称为嵌入超曲面。空集是任意维数的嵌入子流形。

最容易理解的嵌入子流形是余维数为 $0$ 的那些。回忆在\cref{example:ch01-open-submanifolds} 中,对任一光滑流形 $M$,我们把 $M$ 的\textbf{开子流形}定义为任一带有子空间拓扑、并通过限制 $M$ 的坐标卡而得到光滑坐标卡的开子集。

\begin{proposition}[Open Submanifolds]\label{proposition:ch05-open-submanifolds}
设 $M$ 是光滑流形。$M$ 中余维数为 $0$ 的嵌入子流形恰好就是开子流形。
\end{proposition}

\begin{proof}
设 $U\subseteq M$ 是开子流形,并令 $\iota:U\hookrightarrow M$ 为包含映射。\cref{example:ch01-open-submanifolds} 表明 $U$ 是与 $M$ 同维的光滑流形,故其余维数为 $0$。用\cref{example:ch01-open-submanifolds} 中构造的 $U$ 的光滑坐标卡, $\iota$ 在坐标下由恒等映射表示,故它是光滑浸入;又因 $U$ 具有子空间拓扑, $\iota$ 是光滑嵌入。于是 $U$ 是嵌入子流形。反之,设 $U$ 是 $M$ 的任一余维数为 $0$ 的嵌入子流形。则包含映射 $\iota:U\hookrightarrow M$ 由定义是光滑嵌入,因而由\cref{proposition:ch04-local-diffeo-equiv} 是局部微分同胚,并由\cref{proposition:ch04-local-diffeo-properties} 是开映射。故 $U$ 是 $M$ 的开子集。
\end{proof}

下面几个命题展示了生成嵌入子流形的另外几种方式。

\begin{proposition}[Images of Embeddings as Submanifolds]\label{proposition:ch05-images-of-embeddings}
设 $M$ 是带边或不带边的光滑流形, $N$ 是光滑流形, $F:N\to M$ 是光滑嵌入。令 $S=F(N)$。在子空间拓扑下, $S$ 是拓扑流形,并且它具有唯一的光滑结构,使它成为 $M$ 的嵌入子流形且 $F$ 是到其像上的微分同胚。
\end{proposition}

\begin{proof}
若给 $S$ 赋予它从 $M$ 继承来的子空间拓扑,则 $F$ 是嵌入这一假设意味着 $F$ 可视为从 $N$ 到 $S$ 上的同胚,从而 $S$ 是拓扑流形。我们把形如 $\bigl(F(U),\varphi\circ F^{-1}\bigr)$ 的坐标卡取为 $S$ 的光滑坐标卡,其中 $(U,\varphi)$ 是 $N$ 的任一光滑坐标卡;这些坐标卡的光滑相容性直接来自 $N$ 相应坐标卡的光滑相容性。在 $S$ 的这一光滑结构下,映射 $F$ 是到其像上的微分同胚(本质上由定义即如此),而这显然是具有该性质的唯一光滑结构。包含映射 $S\hookrightarrow M$ 等于一个微分同胚接着一个光滑嵌入的复合
\begin{align}\label{align:ch05-inclusion-factor}
S\xrightarrow{\ F^{-1}\ }N\xrightarrow{\ F\ }M,
\end{align}
故它是光滑嵌入。
\end{proof}

由于每个嵌入子流形都是某个光滑嵌入(即它自身的包含映射)的像,上一命题表明嵌入子流形恰好就是光滑嵌入的像。

\begin{proposition}[Slices of Product Manifolds]\label{proposition:ch05-product-slices}
设 $M$ 与 $N$ 为光滑流形。对每个 $p\in N$,子集 $M\times\{p\}$(称为积流形的一个切片)是 $M\times N$ 的嵌入子流形,微分同胚于 $M$。
\end{proposition}

\begin{proof}
集合 $M\times\{p\}$ 是光滑嵌入 $x\mapsto(x,p)$ 的像。
\end{proof}

\begin{proposition}[Graphs as Submanifolds]\label{proposition:ch05-graphs}
设 $M$ 是光滑 $m$ 流形(不带边), $N$ 是带边或不带边的光滑 $n$ 流形, $U\subseteq M$ 开, $f:U\to N$ 是光滑映射。令 $\Gamma(f)\subseteq M\times N$ 表示 $f$ 的图:
\begin{align*}
\Gamma(f)=\{(x,y)\in M\times N:x\in U,\ y=f(x)\}.
\end{align*}
则 $\Gamma(f)$ 是 $M\times N$ 的 $m$ 维嵌入子流形(见\cref{figure:ch05-graph})。
\end{proposition}

\begin{proof}
定义映射 $\gamma_f:U\to M\times N$ 为
\begin{align}\label{align:ch05-graph-map}
\gamma_f(x)=\bigl(x,f(x)\bigr).
\end{align}
它是光滑映射,其像为 $\Gamma(f)$。由于投影 $\pi_M:M\times N\to M$ 满足对 $x\in U$ 有 $\pi_M\circ\gamma_f(x)=x$,故复合 $\mathrm{d}(\pi_M)_{(x,f(x))}\circ\mathrm{d}(\gamma_f)_x$ 对每个 $x\in U$ 是 $T_xM$ 上的恒等映射。因此 $\mathrm{d}(\gamma_f)_x$ 单射,故 $\gamma_f$ 是光滑浸入。它是到其像上的同胚,因为 $\pi_M|_{\Gamma(f)}$ 是它的连续逆。于是 $\Gamma(f)$ 是微分同胚于 $U$ 的嵌入子流形。
\end{proof}

\begin{figure}[ht]
\centering
\includegraphics[width=0.8\textwidth]{fig-05-1.png}
\caption{图像是嵌入子流形}\label{figure:ch05-graph}
\end{figure}

对某些用途而言,仅仅是嵌入子流形这一条件还不够强。(例如,见后面的\cref{lemma:ch05-extension}。)嵌入子流形 $S\subseteq M$ 称为\textbf{真嵌入的},如果包含映射 $S\hookrightarrow M$ 是固有映射。

\begin{proposition}\label{proposition:ch05-proper-closed}
设 $M$ 是带边或不带边的光滑流形, $S\subseteq M$ 是嵌入子流形。则 $S$ 真嵌入,当且仅当它是 $M$ 的闭子集。
\end{proposition}

\begin{proof}
若 $S$ 真嵌入,则由\cref{theorem:appA-proper-closed} 知它闭。反之,若 $S$ 在 $M$ 中闭,则\cref{proposition:appA-proper-criterion}(c) 表明包含映射 $S\hookrightarrow M$ 固有。
\end{proof}

\begin{corollary}\label{corollary:ch05-compact-properly-embedded}
每个紧嵌入子流形都是真嵌入的。
\end{corollary}

\begin{proof}
Hausdorff 空间中的紧子集是闭的。
\end{proof}

整体定义函数的图是真嵌入子流形的常见例子。

\begin{proposition}[Global Graphs Are Properly Embedded]\label{proposition:ch05-global-graph-proper}
设 $M$ 是光滑流形, $N$ 是带边或不带边的光滑流形, $f:M\to N$ 是光滑映射。在\cref{proposition:ch05-graphs} 的光滑流形结构下, $\Gamma(f)$ 在 $M\times N$ 中真嵌入。
\end{proposition}

\begin{proof}
在这种情形下,投影 $\pi_M:M\times N\to M$ 是由\eqref{align:ch05-graph-map} 定义的嵌入 $\gamma_f:M\to M\times N$ 的光滑左逆。故由\cref{proposition:appA-proper-criterion} 知 $\gamma_f$ 固有。
\end{proof}

\subsection{嵌入子流形的切片坐标卡}\label{subsection:ch05-slice-charts}

正如我们下一定理将表明的,嵌入子流形局部地以 $\mathbb{R}^k$ 到 $\mathbb{R}^n$ 的标准嵌入为模型,即把 $\mathbb{R}^k$ 等同于子空间
\begin{align*}
\bigl\{\bigl(x^1,\dots,x^k,x^{k+1},\dots,x^n\bigr):x^{k+1}=\cdots=x^n=0\bigr\}\subseteq\mathbb{R}^n.
\end{align*}
更一般地,若 $U$ 是 $\mathbb{R}^n$ 的开子集且 $k\in\{0,\dots,n\}$, $U$ 的一个 $k$ 维\textbf{切片}(或简称 $k$-切片)是指形如
\begin{align*}
S=\bigl\{\bigl(x^1,\dots,x^k,x^{k+1},\dots,x^n\bigr)\in U:x^{k+1}=c^{k+1},\dots,x^n=c^n\bigr\}
\end{align*}
的任一子集,其中 $c^{k+1},\dots,c^n$ 为常数。(当 $k=n$ 时,这恰好意味着 $S=U$。)显然,每个 $k$-切片都同胚于 $\mathbb{R}^k$ 的开子集。(有时把除最后几个坐标外的某个坐标子集取为常数来定义切片是方便的。其含义应可由上下文清楚判断。)

设 $M$ 是光滑 $n$ 流形, $(U,\varphi)$ 是 $M$ 上的光滑坐标卡。若 $S$ 是 $U$ 的子集且 $\varphi(S)$ 是 $\varphi(U)$ 的 $k$-切片,则称 $S$ 是 $U$ 的 $k$-切片。(虽然一般我们允许切片由任意常数 $c^{k+1},\dots,c^n$ 定义,但有时用常数全为零的切片坐标是方便的,这只需从每个坐标函数减去一个常数即可轻松实现。)给定子集 $S\subseteq M$ 与非负整数 $k$,我们称 $S$ 满足\textbf{局部 $k$-切片条件},如果 $S$ 的每一点都含于 $M$ 的某光滑坐标卡 $(U,\varphi)$ 的定义域内,使 $S\cap U$ 是 $U$ 中单个 $k$-切片。任何这样的坐标卡称为 $S$ 在 $M$ 中的\textbf{切片坐标卡},相应的坐标 $\bigl(x^1,\dots,x^n\bigr)$ 称为\textbf{切片坐标}。

\begin{theorem}[Local Slice Criterion for Embedded Submanifolds]\label{theorem:ch05-slice-criterion}
设 $M$ 是光滑 $n$ 流形。若 $S\subseteq M$ 是 $k$ 维嵌入子流形,则 $S$ 满足局部 $k$-切片条件。反之,若 $S\subseteq M$ 是满足局部 $k$-切片条件的子集,则在子空间拓扑下, $S$ 是 $k$ 维拓扑流形,并具有一个光滑结构使它成为 $M$ 的 $k$ 维嵌入子流形。
\end{theorem}

\begin{proof}
先设 $S\subseteq M$ 是 $k$ 维嵌入子流形。由于包含映射 $S\hookrightarrow M$ 是浸入,秩定理表明对任一 $p\in S$,存在 $S$(在其给定光滑流形结构下)的坐标卡 $(U,\varphi)$ 与 $M$ 的坐标卡 $(V,\psi)$,二者都以 $p$ 为中心,使包含映射 $\iota|_U:U\to V$ 具有坐标表示
\begin{align*}
\bigl(x^1,\dots,x^k\bigr)\mapsto\bigl(x^1,\dots,x^k,0,\dots,0\bigr).
\end{align*}
取 $\varepsilon>0$ 充分小,使 $U$ 与 $V$ 都含以 $p$ 为中心、半径为 $\varepsilon$ 的坐标球,并记这些坐标球为 $U_0\subseteq U$ 与 $V_0\subseteq V$。于是 $\iota(U_0)$ 恰好是 $V_0$ 中的单个切片。由于 $S$ 具有子空间拓扑, $U_0$ 在 $S$ 中开这一事实意味着存在开子集 $W\subseteq M$ 使 $U_0=W\cap S$。令 $V_1=V_0\cap W$,我们得到 $M$ 的一个含 $p$ 的光滑坐标卡 $(V_1,\psi|_{V_1})$,使 $V_1\cap S=U_0$,而这是 $V_1$ 的单个切片。

反之,设 $S$ 满足局部 $k$-切片条件。在子空间拓扑下, $S$ Hausdorff 且第二可数,因为这两条性质都被子空间继承。为看 $S$ 局部 Euclid,我们构造一个图册。构造的基本想法是:若 $\bigl(x^1,\dots,x^n\bigr)$ 是 $S$ 在 $M$ 中的切片坐标,则我们可用 $\bigl(x^1,\dots,x^k\bigr)$ 作为 $S$ 的局部坐标。

在本证明中,令 $\pi:\mathbb{R}^n\to\mathbb{R}^k$ 表示投向前 $k$ 个坐标的投影。任取 $S$ 在 $M$ 中的切片坐标卡 $(U,\varphi)$(\cref{figure:ch05-slice-chart}),并定义
\begin{align*}
V=U\cap S,\qquad \widehat{V}=\pi\circ\varphi(V),\qquad \psi=\pi\circ\varphi|_V:V\to\widehat{V}.
\end{align*}
由切片坐标卡的定义, $\varphi(V)$ 是 $\varphi(U)$ 与某 $k$-切片 $A\subseteq\mathbb{R}^n$ 的交,其中 $A$ 由令 $x^{k+1}=c^{k+1},\dots,x^n=c^n$ 定义,故 $\varphi(V)$ 在 $A$ 中开。由于 $\pi|_A$ 是从 $A$ 到 $\mathbb{R}^k$ 的微分同胚, $\widehat{V}$ 在 $\mathbb{R}^k$ 中开。此外, $\psi$ 是同胚,因为它有连续逆 $\varphi^{-1}\circ j|_{\widehat{V}}$,其中 $j:\mathbb{R}^k\to\mathbb{R}^n$ 是映射
\begin{align*}
j\bigl(x^1,\dots,x^k\bigr)=\bigl(x^1,\dots,x^k,c^{k+1},\dots,c^n\bigr).
\end{align*}
于是 $S$ 是拓扑 $k$ 流形,且包含映射 $\iota:S\hookrightarrow M$ 是拓扑嵌入。

为在 $S$ 上赋予光滑结构,我们需验证上面构造的坐标卡光滑相容。设 $(U,\varphi)$ 与 $(U',\varphi')$ 是 $S$ 在 $M$ 中的两个切片坐标卡,并设 $(V,\psi)$、$(V',\psi')$ 是 $S$ 的相应坐标卡。转移映射由 $\psi'\circ\psi^{-1}=\pi\circ\varphi'\circ\varphi^{-1}\circ j$ 给出,它是四个光滑映射的复合(\cref{figure:ch05-slice-compat})。于是我们构造的图册实际上是光滑图册,它在 $S$ 上定义了一个光滑结构。用 $M$ 的切片坐标卡 $(U,\varphi)$ 与 $S$ 的相应坐标卡 $(V,\psi)$ 表示,包含映射 $S\hookrightarrow M$ 具有坐标表示
\begin{align*}
\bigl(x^1,\dots,x^k\bigr)\mapsto\bigl(x^1,\dots,x^k,c^{k+1},\dots,c^n\bigr),
\end{align*}
这是光滑浸入。由于包含映射既是光滑浸入又是拓扑嵌入,故 $S$ 是嵌入子流形。
\end{proof}

\begin{figure}[ht]
\centering
\includegraphics[width=0.8\textwidth]{fig-05-2.png}
\caption{满足 $k$-切片条件的子集的图}\label{figure:ch05-slice-chart}
\end{figure}

\begin{figure}[ht]
\centering
\includegraphics[width=0.8\textwidth]{fig-05-3.png}
\caption{切片图的光滑相容性}\label{figure:ch05-slice-compat}
\end{figure}

注意, $S\subseteq M$ 的局部切片条件只是关于子集 $S$ 本身的条件;它不以 $S$ 上任何特定拓扑或光滑结构为前提。正如我们后面将看到的(\cref{theorem:ch05-uniqueness}),上一定理中构造的光滑流形结构是使 $S$ 可被视为子流形的唯一结构,故满足局部切片条件的子集只以一种方式成为嵌入子流形。

\begin{example}[Spheres as Submanifolds]\label{example:ch05-spheres-slice}
对任意 $n\geqslant 0$, $S^n$ 是 $\mathbb{R}^{n+1}$ 的嵌入子流形,因为它局部是光滑函数的图:正如我们在\cref{example:ch01-sphere-graph} 中证明的, $S^n$ 与开子集 $\{x:x^i>0\}$ 的交是光滑函数
\begin{align*}
x^i=f\bigl(x^1,\dots,x^{i-1},x^{i+1},\dots,x^{n+1}\bigr)
\end{align*}
的图,其中 $f:B^n\to\mathbb{R}$ 由 $f(u)=\sqrt{1-|u|^2}$ 给出。类似地, $S^n$ 与 $\{x:x^i<0\}$ 的交是 $-f$ 的图。由于 $S^n$ 中每一点都在这些集合之一中, $S^n$ 满足局部 $n$-切片条件,从而是 $\mathbb{R}^{n+1}$ 的嵌入子流形。如此诱导到 $S^n$ 上的光滑结构与我们在第一章中定义的相同:事实上,这些切片坐标卡所确定的 $S^n$ 坐标恰好就是\cref{example:ch01-graph-coordinates} 中定义的图坐标。
\end{example}

\begin{exercise}\label{exercise:ch05-spherical-coords}
证明球坐标(\cref{example:appC-spherical-coords})在 $\mathbb{R}^3$ 中任一其有定义的开子集上构成 $S^2$ 的切片坐标卡。
\end{exercise}

若 $M$ 是边界非空的光滑流形且 $S\subseteq M$ 是嵌入子流形,则 $S$ 可能以非常复杂的方式与 $\partial M$ 相交,因此在那种情形下我们不试图证明关于 $S$ 在 $M$ 中切片坐标卡存在性的任何一般结果。然而,在子流形恰为 $M$ 自身边界这一特殊情形下, $M$ 的边界坐标卡就扮演了 $S$ 在 $M$ 中切片坐标卡的角色;而我们确实有如下结果。

\begin{theorem}\label{theorem:ch05-boundary-submanifold}
若 $M$ 是带边光滑 $n$ 流形,则在子空间拓扑下, $\partial M$ 是 $(n-1)$ 维拓扑流形(不带边),并具有一个光滑结构使它成为 $M$ 的真嵌入子流形。
\end{theorem}

\begin{proof}
见习题 5-2。
\end{proof}

我们后面将看到 $\partial M$ 上的光滑结构是唯一的。(见\cref{theorem:ch05-uniqueness}。)为了在 $M$ 有边界时分析 $M$ 的更一般子流形,最有效的技巧通常是把 $M$ 嵌入一个更大的不带边光滑流形 $\widehat{M}$ 中,并在该背景下应用前面的结果。\cref{example:ch09-embedding-boundary} 将表明每个带边光滑流形都能以这种方式嵌入。

\subsection{水平集}\label{subsection:ch05-level-sets}

实际上,嵌入子流形最常以方程或方程组的解集形式出现。推广我们在\cref{example:ch01-level-sets} 中引进的术语,若 $\Phi:M\to N$ 是任一映射, $c$ 是 $N$ 的任一点,则称集合 $\Phi^{-1}(c)$ 为 $\Phi$ 的一个\textbf{水平集}(\cref{figure:ch05-level-set})。(在 $N=\mathbb{R}^k$ 且 $c=0$ 的特殊情形,水平集 $\Phi^{-1}(0)$ 通常称为 $\Phi$ 的零点集。)

\begin{figure}[ht]
\centering
\includegraphics[width=0.8\textwidth]{fig-05-4.png}
\caption{水平集}\label{figure:ch05-level-set}
\end{figure}

容易找到不是光滑子流形的光滑函数水平集。例如,考察由
\begin{align*}
\rho(x,y)=x^2-y,\qquad \Phi(x,y)=x^2-y^2,\qquad \Psi(x,y)=x^2-y^3
\end{align*}
定义的三个光滑函数 $\rho,\Phi,\Psi:\mathbb{R}^2\to\mathbb{R}$(见\cref{figure:ch05-bad-level-sets})。虽然 $\rho$ 的零点集(一条抛物线)是 $\mathbb{R}^2$ 的嵌入子流形,因为它是光滑函数 $f(x)=x^2$ 的图,但习题 5-11 要求你证明 $\Phi$ 与 $\Psi$ 的零点集都不是嵌入子流形。事实上,若不对光滑函数作进一步假设,情况差不多坏到可以想象的地步:正如\cref{theorem:ch01-closed-zero-set} 所表明的, $M$ 的每个闭子集都可表为某光滑实值函数的零点集。

\begin{figure}[ht]
\centering
\includegraphics[width=0.8\textwidth]{fig-05-5.png}
\caption{水平集可能是也可能不是嵌入子流形}\label{figure:ch05-bad-level-sets}
\end{figure}

我们在\cref{example:ch01-level-sets} 中使用的论证(基于隐函数定理)用以证明 $\mathbb{R}^n$ 中某些水平集是光滑流形,可加以改造以证明那些水平集实际上是 $\mathbb{R}^n$ 的嵌入子流形。但利用秩定理,我们能证明强得多的结论。

\begin{theorem}[Constant-Rank Level Set Theorem]\label{theorem:ch05-constant-rank-level-set}
设 $M$ 与 $N$ 为光滑流形, $\Phi:M\to N$ 是秩为常数 $r$ 的光滑映射。则 $\Phi$ 的每个水平集都是 $M$ 中余维数为 $r$ 的真嵌入子流形。
\end{theorem}

\begin{proof}
记 $m=\dim M$, $n=\dim N$, $k=m-r$。任取 $c\in N$,并令 $S$ 表示水平集 $\Phi^{-1}(c)\subseteq M$。由秩定理,对每个 $p\in S$,存在以 $p$ 为中心的坐标卡 $(U,\varphi)$ 与以 $c=\Phi(p)$ 为中心的坐标卡 $(V,\psi)$,使 $\Phi$ 具有形如\eqref{align:ch04-rank-canonical} 的坐标表示,因而 $S\cap U$ 是切片
\begin{align*}
\bigl\{\bigl(x^1,\dots,x^r,x^{r+1},\dots,x^m\bigr)\in U:x^1=\cdots=x^r=0\bigr\}.
\end{align*}
于是 $S$ 满足局部 $k$-切片条件,故它是 $k$ 维嵌入子流形。它由连续性在 $M$ 中闭,故由\cref{proposition:ch05-proper-closed} 知它真嵌入。
\end{proof}

\begin{corollary}[Submersion Level Set Theorem]\label{corollary:ch05-submersion-level-set}
若 $M$ 与 $N$ 为光滑流形, $\Phi:M\to N$ 是光滑淹没,则 $\Phi$ 的每个水平集都是真嵌入子流形,其余维数等于 $N$ 的维数。
\end{corollary}

\begin{proof}
每个光滑淹没都具有等于其陪域维数的常秩。
\end{proof}

这一结果应与线性代数中的相应结果相比较:若 $L:\mathbb{R}^m\to\mathbb{R}^r$ 是满射线性映射,则由秩-零维律, $L$ 的核是余维数为 $r$ 的线性子空间。向量方程 $Lx=0$ 等价于 $r$ 个线性无关的数量方程,每个都可视为削减 $\mathbb{R}^m$ 中的一个自由度,留下余维数为 $r$ 的子空间。在光滑流形的背景下,满射线性映射的类比是光滑淹没,其每个(局部)分量函数都把维数削减一。

\cref{corollary:ch05-submersion-level-set} 可大大加强,因为我们只需在所关心的水平集上检验淹没条件。若 $\Phi:M\to N$ 是光滑映射,称点 $p\in M$ 是 $\Phi$ 的\textbf{正则点},如果 $\mathrm{d}\Phi_p:T_pM\to T_{\Phi(p)}N$ 满射;否则称它是 $\Phi$ 的\textbf{临界点}。这特别意味着:若 $\dim M<\dim N$,则 $M$ 的每一点都是临界点;而每一点正则当且仅当 $F$ 是淹没。注意,由\cref{proposition:ch04-full-rank-open}, $\Phi$ 的正则点集总是 $M$ 的开子集。称点 $c\in N$ 是 $\Phi$ 的\textbf{正则值},如果水平集 $\Phi^{-1}(c)$ 的每一点都是正则点;否则称它是\textbf{临界值}。特别地,若 $\Phi^{-1}(c)=\varnothing$,则 $c$ 是正则值。最后,若 $c$ 是 $\Phi$ 的正则值,则称水平集 $\Phi^{-1}(c)$ 为\textbf{正则水平集};换言之,正则水平集是完全由 $\Phi$ 的正则点(即 $\mathrm{d}\Phi_p$ 满射的点 $p$)组成的水平集。

\begin{corollary}[Regular Level Set Theorem]\label{corollary:ch05-regular-level-set}
光滑流形之间光滑映射的每个正则水平集都是真嵌入子流形,其余维数等于陪域的维数。
\end{corollary}

\begin{proof}
设 $\Phi:M\to N$ 是光滑映射, $c\in N$ 是正则值。由\cref{proposition:ch04-full-rank-open},使 $\operatorname{rank}\mathrm{d}\Phi_p=\dim N$ 的点 $p\in M$ 之集 $U$ 在 $M$ 中开,并且由 $c$ 为正则值的假设它包含 $\Phi^{-1}(c)$。于是 $\Phi|_U:U\to N$ 是光滑淹没,而上一推论表明 $\Phi^{-1}(c)$ 是 $U$ 的嵌入子流形。由于光滑嵌入之复合 $\Phi^{-1}(c)\hookrightarrow U\hookrightarrow M$ 仍是光滑嵌入,故 $\Phi^{-1}(c)$ 是 $M$ 的嵌入子流形;并且它由连续性闭。
\end{proof}

值得注意,上一推论也适用于空水平集,它既是正则水平集又是真嵌入子流形。

\begin{example}[Spheres]\label{example:ch05-sphere-regular-level}
现在我们可以给出 $S^n$ 是 $\mathbb{R}^{n+1}$ 的嵌入子流形的一个简单得多的证明。球面是光滑函数 $f:\mathbb{R}^{n+1}\to\mathbb{R}$(由 $f(x)=|x|^2$ 给出)的正则水平集,因为 $\mathrm{d}f_x(v)=2\sum_i x^iv_i$,它除在原点外都满射。
\end{example}

并非每个嵌入子流形都能表为光滑淹没的水平集。然而,下一个命题表明每个嵌入子流形至少局部地具有这种形式。

\begin{proposition}\label{proposition:ch05-local-level-set}
设 $S$ 是光滑 $m$ 流形 $M$ 的子集。则 $S$ 是 $M$ 的嵌入 $k$-子流形,当且仅当 $S$ 的每一点都有 $M$ 中的邻域 $U$,使 $U\cap S$ 是某光滑淹没 $\Phi:U\to\mathbb{R}^{m-k}$ 的水平集。
\end{proposition}

\begin{proof}
先设 $S$ 是嵌入 $k$-子流形。若 $\bigl(x^1,\dots,x^m\bigr)$ 是 $S$ 在开子集 $U\subseteq M$ 上的切片坐标,则由 $\Phi(x)=\bigl(x^{k+1},\dots,x^m\bigr)$ 在坐标下给出的映射 $\Phi:U\to\mathbb{R}^{m-k}$ 容易看出是光滑淹没,其一个水平集为 $S\cap U$(\cref{figure:ch05-local-level-set})。反之,设围绕每一点 $p\in S$ 都有邻域 $U$ 与光滑淹没 $\Phi:U\to\mathbb{R}^{m-k}$ 使 $S\cap U$ 是 $\Phi$ 的水平集。由淹没水平集定理, $S\cap U$ 是 $U$ 的嵌入子流形,故它满足局部切片条件;由此即得 $S$ 本身是 $M$ 的嵌入子流形。
\end{proof}

\begin{figure}[ht]
\centering
\includegraphics[width=0.8\textwidth]{fig-05-6.png}
\caption{嵌入子流形局部是水平集}\label{figure:ch05-local-level-set}
\end{figure}

若 $S\subseteq M$ 是嵌入子流形,使 $S$ 为 $\Phi$ 的正则水平集的光滑映射 $\Phi:M\to N$ 称为 $S$ 的\textbf{定义映射}。在 $N=\mathbb{R}^{m-k}$(即 $\Phi$ 是实值或向量值函数)的特殊情形,它通常称为\textbf{定义函数}。\cref{example:ch05-sphere-regular-level} 表明 $f(x)=|x|^2$ 是球面的定义函数。更一般地,若 $U$ 是 $M$ 的开子集, $\Phi:U\to N$ 是使 $S\cap U$ 为 $\Phi$ 的正则水平集的光滑映射,则 $\Phi$ 称为 $S$ 的\textbf{局部定义映射}(或局部定义函数)。\cref{proposition:ch05-local-level-set} 说每个嵌入子流形在其每一点的邻域内都容许局部定义函数。

在具体例子中,为子流形寻找(局部或整体)定义函数通常只需结合关于子流形如何定义的几何信息与一些计算上的巧思。下面是一个例子。

\begin{example}[Surfaces of Revolution]\label{example:ch05-surface-revolution}
令 $H$ 为半平面 $\{(r,z):r>0\}$,并设 $C\subseteq H$ 是 $1$ 维嵌入子流形。由 $C$ 确定的旋转曲面是 $\mathbb{R}^3$ 的子集 $S_C$,由
\begin{align*}
S_C=\left\{(x,y,z):\bigl(\sqrt{x^2+y^2},z\bigr)\in C\right\}
\end{align*}
给出。集合 $C$ 称为它的\textbf{生成曲线}(见\cref{figure:ch05-surface-revolution})。若 $\varphi:U\to\mathbb{R}$ 是 $C$ 在 $H$ 中的任一局部定义函数,则我们由
\begin{align*}
\Phi(x,y,z)=\varphi\bigl(\sqrt{x^2+y^2},z\bigr)
\end{align*}
得到 $S_C$ 的局部定义函数 $\Phi$,它定义在开子集
\begin{align*}
\widetilde{U}=\left\{(x,y,z):\bigl(\sqrt{x^2+y^2},z\bigr)\in U\right\}\subseteq\mathbb{R}^3
\end{align*}
上。计算表明 $\Phi$ 的 Jacobian 矩阵为
\begin{align*}
\mathrm{D}\Phi(x,y,z)=\begin{pmatrix}
\dfrac{x}{r}\dfrac{\partial\varphi}{\partial r}(r,z) & \dfrac{y}{r}\dfrac{\partial\varphi}{\partial r}(r,z) & \dfrac{\partial\varphi}{\partial z}(r,z)
\end{pmatrix},
\end{align*}
其中我们记 $r=\sqrt{x^2+y^2}$。在任一点 $(x,y,z)\in S_C$ 处, $\mathrm{D}\Phi(x,y,z)$ 的分量中至少一个非零,故 $S_C$ 是 $\Phi$ 的正则水平集,从而是 $\mathbb{R}^3$ 的 $2$ 维嵌入子流形。

作为一个具体例子, \cref{example:ch04-submersions-immersions}(d) 中描述的甜甜圈状旋转环面 $D$ 是由圆 $(r-2)^2+z^2=1$ 得到的旋转曲面。它是函数 $\Phi(x,y,z)=\bigl(\sqrt{x^2+y^2}-2\bigr)^2+z^2$ 的正则水平集,该函数在 $\mathbb{R}^3$ 除去 $z$ 轴外光滑。
\end{example}

\begin{figure}[ht]
\centering
\includegraphics[width=0.8\textwidth]{fig-05-7.png}
\caption{旋转曲面}\label{figure:ch05-surface-revolution}
\end{figure}

\end{document}